%! Equations and Graphs in Both Directions — Unit 2, Lesson 2.3 — Lesson Plan
\documentclass[10pt]{article}
\usepackage{atda-article}
\usepackage{atda-boxes}
\usepackage{graphicx}   % -article does NOT load graphicx; needed for thumbnails/screenshots
\graphicspath{{images/}}
\newcommand{\TallMath}[1]{$\displaystyle #1\rule[-1.4em]{0pt}{3.2em}$}
\newcommand{\UnitNumberName}{Unit 2: Functions, Transformations, and Graph Analysis \quad}
\newcommand{\LessonNumberName}{Lesson 2.3: Equations and Graphs in Both Directions}

\begin{document}
\begin{center}
    {\Large\bfseries \CourseName \\
    {\small\bfseries \UnitNumberName \LessonNumberName}}
\end{center}

\begin{tcolorbox}[colback=mist, colframe=royal, boxrule=0.9pt, arc=2mm,
    left=2mm, right=2mm, top=1.5mm, bottom=1.5mm]
    \textbf{Primary Objective:} I can graph a linear, quadratic, or exponential function from its
    equation by moving the parent's anchor point, and I can write the equation of a function from its
    graph by reading the vertex, the asymptote, or two points. When the shape is not the parent's, I
    can find the stretch factor $a$ from one extra point. I can use a graphing calculator to verify a
    transformation --- typing it correctly and setting a window that actually shows it --- and I can
    explain why a matching screen is not proof and a substitution is.
\end{tcolorbox}

\renewcommand{\arraystretch}{1.6}
\begin{skillbox}[Priority Ideas \& Skills]{goldbox}
\begin{tabularx}{\linewidth}{@{}X|X@{}}
\begin{minipage}[t]{\linewidth}\vspace{0pt}
\textbf{Standards covered:}
\begin{itemize}[leftmargin=1.1em, itemsep=2pt, topsep=2pt]
  \item \texttt{AFDA.AF.1d} --- write the equation of a linear, quadratic, or exponential function
        \emph{given a graph}, using transformations of the parent function.
  \item \texttt{AFDA.AF.1f} --- graph a function \emph{given its equation}, using transformations of
        the parent function. \emph{Use technology to verify transformations of functions.}
\end{itemize}
\textbf{Skills students leave with:}
\begin{itemize}[leftmargin=1.1em, itemsep=2pt, topsep=2pt]
  \item Locate the \textbf{anchor point} of an image from its equation, and carry four more points
        with it onto a pre-drawn grid.
  \item Read a graph's landmark --- \textbf{vertex}, \textbf{asymptote}, or \textbf{two points} ---
        and write the equation from it.
  \item Solve for $a$ in $y=a(x-h)^2+k$ from one additional point, and say what $|a|$ does to the
        shape.
  \item Type an equation into a TI-84 correctly (parentheses, the \texttt{(-)} key) and set a
        \textbf{window} from the anchor point without graphing first.
  \item \textbf{Verify by substitution}, and explain why a matching screen is not proof.
\end{itemize}
\end{minipage}
&
\begin{minipage}[t]{\linewidth}\vspace{0pt}
\textbf{Key Understandings (the \emph{why}):}
\begin{itemize}[leftmargin=1.1em, itemsep=2pt, topsep=2pt]
  \item \textbf{2.2 asked them to \emph{describe} a move they were shown. 2.3 asks them to
        \emph{produce} the other representation.} That is the whole difference, and it is why
        \texttt{1d} and \texttt{1f} are one lesson: they are the same translation run in opposite
        directions.
  \item \textbf{An anchor point is the entire method.} The AFDA guidance calls the parent graph an
        ``anchor graph from which other graphs are derived''; one tracked point turns a
        transformation into arithmetic on coordinates.
  \item \textbf{``Use technology to \emph{verify}'' is the standard's own word, and verify means
        check --- not decide.} A screen is a picture at pixel resolution; a substitution is a number.
        When they disagree, the number wins and there is a typo to find.
  \item \textbf{The parentheses question from 2.2 is now literally a keystroke.}
        \texttt{2\^{}X-5} versus \texttt{2\^{}(X-5)} is outside-versus-inside with a calculator
        attached, which is why the hook is a typing error.
  \item \textbf{The guidance singles out lines: ``the equation of a linear function can be determined
        by two points\dots or by the slope and a point.''} It says that because for a line the
        transformations are not distinguishable --- $2f(x)$ and $f(2x)$ are the same graph. Teach the
        two-point method there and say why.
\end{itemize}
\end{minipage}
\end{tabularx}
\end{skillbox}

\begin{skillbox}[{Vocabulary, Concepts, \& Theorems}]{greenbox}
\renewcommand{\arraystretch}{1.35}
\begin{tabularx}{\linewidth}{@{}p{3.6cm}X@{}}
\textbf{Term} & \textbf{Meaning students record} \\
\hline
Anchor point & The one point on the parent you track through a transformation: the turning point of
    a parabola, the point $(0,1)$ on an exponential. Where it lands gives both translations. \\
Pre-image and image & The graph before the transformation is the \emph{pre-image}; the graph after
    it is the \emph{image}. (The AFDA guidance's own words.) \\
Transformation form & An equation written so every move is visible: $y=a\,f(x-h)+k$ for parent $f$.
    For a quadratic, \TallMath{y=a(x-h)^2+k} with vertex $(h,k)$. \\
Viewing window & The rectangle of the plane the calculator screen shows, set by \texttt{Xmin},
    \texttt{Xmax}, \texttt{Ymin}, \texttt{Ymax}. A correct graph can be entirely off-screen. \\
Substitution check & Putting one input into your equation and comparing the output with the graph or
    the \textsc{table}. The only thing that proves an equation matches a picture. \\
Stretch factor $a$ & The number multiplying the parent. $|a|>1$ narrower and steeper; $0<|a|<1$
    flatter and wider; $a<0$ also flips it. \\
\end{tabularx}
\end{skillbox}

\begin{fixedskillbox}[Activate Prior Knowledge \& Spiral Review]{mist}
\begin{tabularx}{\linewidth}{@{}X c@{}}
\begin{minipage}[t]{0.55\linewidth}\vspace{0pt}
\textbf{The warm-up is five items, and item 2 is the hook planted:}
\begin{enumerate}[leftmargin=1.4em, itemsep=1pt, topsep=2pt]
  \item A table for $f(x)=(x-2)^2$, plus which input gives the smallest output --- feeds notes box 1
  \item $2^x-5$ versus $2^{x-5}$ at $x=5$ --- the whole hook, before it is asked
  \item Name the parent and the transformation (Lesson 2.2 spiral)
  \item Domain and range in interval notation (Lesson 2.1 spiral)
  \item Last night's item 10: what did you predict the calculator would show?
\end{enumerate}
\end{minipage}
&
\begin{minipage}[t]{0.38\linewidth}\vspace{0pt}
\textbf{How to use the results:} item 2 tells you who can already keep a number inside the
parentheses --- $27$ against $1$ from the same input. Item 3(c) is the negative-key trap in advance.
Item 5 is homework accountability: if nobody can answer it, box 1 will need the time you budgeted
for box 3.
\end{minipage}
\end{tabularx}
\end{fixedskillbox}

\begin{skillbox}[Hook]{mist}
\textbf{Two students, one instruction, two different screens.} Put both keystroke strings on the
board: \texttt{Y1=2\^{}X-5} and \texttt{Y1=2\^{}(X-5)}. The instruction they were both given was
\emph{graph the exponential parent moved $5$ units to the right}. \textbf{Pose it this way:} ``Both of
them got a curve on the screen. Only one of them followed the instruction. How do you decide which,
without asking me?'' \textbf{Do not take a vote this time} --- 2.2's hook already spent the vote, and
today's point is the opposite one: you do not need consensus, you need one number. Moved right $5$
means the image at $x=5$ must show what the parent showed at $x=0$, namely $2^0=1$. A gives
$2^5-5=27$; B gives $2^{5-5}=1$. \textbf{That is the whole lesson in two lines}, and it is why
\texttt{AF.1f} says \emph{verify}, not \emph{look}.
\end{skillbox}

\boxguard
\begin{skillbox}[Lesson]{mist}
\begin{multicols}{2}
\textbf{1. Equation $\to$ graph (12 min).} Work the routine on $y=(x-2)^2-3$. Fill the four in-text
blanks together, then the parent-to-image table, then \textbf{let them plot} --- the grid is
pre-drawn and the parent is already on it, so nobody sketches from a blank page. \textbf{Ask:} ``What
did the five image points cost you?'' The answer --- adding $2$ and subtracting $3$, five times --- is
the point. Close with where the anchor sits for the other two parents.

\textbf{2. Graph $\to$ equation (14 min).} Three panels, three landmarks. Do the quadratic and the
exponential as a class; assign the line. \textbf{The exponential row is the one to slow down on:}
students look for a vertex and there isn't one. Point at the dashed line and ask what it used to be.
Then the $a$ subsection: vertex gives two numbers, a second point gives the third, a \emph{third}
point checks all three. \textbf{Read the warning about lines aloud} --- it is honest mathematics, and
it is the guidance's own reason for the two-point rule.

\textbf{3. The calculator (12 min).} Resolve the hook on the two panels, then the three traps. If
you have a TI-84 on the projector, type \texttt{2\^{}X-5} live, then add the parentheses and let them
watch the curve jump. \textbf{Trap 3 is worth a demonstration:} type $y=(x-20)^2+900$, press
\texttt{ZOOM}\,\texttt{6}, and let the empty screen sit there for ten seconds before anyone explains
it.

\textbf{4. The rule (4 min).} ``A matching picture is not proof.'' Give them $(x-3)^2$ against
$(x-3)^2+0.5$ and let the class try to find a window that tells them apart. Then substitute.

\textbf{5. Guided practice --- the hoodie (12 min).} Vertex, then $a$ from $(5,0)$, then $P(23)=864$.
\textbf{Ask before they start:} ``Why can't you just read \$23 off the graph?'' The honest answer ---
because it falls between the gridlines, and so does the profit --- is the reason the whole lesson
exists.
\end{multicols}
\end{skillbox}

\boxguard
\begin{skillbox}[Explicit Instruction: The Two Directions, and the Check]{mist}
\begin{tabularx}{\linewidth}{@{}X|X@{}}
\begin{minipage}[t]{\linewidth}\vspace{0pt}
\textbf{Equation $\to$ graph (\texttt{AF.1f}):}
\begin{enumerate}[leftmargin=1.4em, itemsep=1pt, topsep=2pt]
  \item Name the parent.
  \item Move the \textbf{anchor point}, using 2.2's outside/inside rule.
  \item Carry four of the parent's points by the same slide.
  \item Join them; state the domain and range.
\end{enumerate}
\textbf{Graph $\to$ equation (\texttt{AF.1d}):}
\begin{enumerate}[leftmargin=1.4em, itemsep=1pt, topsep=2pt]
  \item Name the family from the shape.
  \item Read the landmark: \textbf{vertex} / \textbf{asymptote and the $y$-axis point} /
        \textbf{two points}.
  \item If the shape is not the parent's, substitute one more point and solve for $a$.
  \item \textbf{Check at a point you have not used yet.}
\end{enumerate}
\end{minipage}
&
\begin{minipage}[t]{\linewidth}\vspace{0pt}
\textbf{The four errors to name before they happen:}
\begin{itemize}[leftmargin=1.1em, itemsep=2pt, topsep=2pt]
  \item \textbf{The parentheses typed wrong.} \texttt{2\^{}X-5} for ``right 5.'' The hook, and Tier E
        item 1(a).
  \item \textbf{\texttt{(-X)\^{}2} for a reflection.} It draws the parent, because $(-x)^2=x^2$ ---
        2.2's warm-up item 3, as keystrokes. Tier E item 1(b).
  \item \textbf{Sign error reading a vertex.} A vertex at $(-3,2)$ gives $(x+3)^2+2$. Have them
        substitute $x=-3$ and confirm they get $2$; nothing else convinces.
  \item \textbf{$a$ read as the $y$-value instead of solved for.} Insist on writing the equation with
        $a$ in it and \emph{then} substituting the point.
\end{itemize}
\textbf{Say this out loud:} ``The calculator draws exactly what you typed. Never what you meant.''
\end{minipage}
\end{tabularx}
\end{skillbox}

\begin{skillbox}[Active Monitoring]{redbox}
\begin{itemize}[leftmargin=1.2em, itemsep=2pt, topsep=2pt]
  \item \textbf{Notes box 1, the plotting:} watch for students plotting the \emph{parent's} $y$-values
        at the \emph{image's} $x$-values --- they applied one slide and forgot the other. The prompt
        is ``read me your anchor point.''
  \item \textbf{Notes box 2, the exponential row:} anyone hunting for a vertex needs redirecting to
        the dashed line. Ask ``where was that line before?''
  \item \textbf{Notes box 2, the $a$ example:} listen for ``$a=3$ because the point is $3$.'' Make
        them write the equation with $a$ in it first.
  \item \textbf{Activity, Tier A item 1:} groups often get $a=1$ by reading the vertex as the second
        point. Two \emph{different} points are needed; say so.
  \item \textbf{Cold-call prompts:} ``Name the parent.'' ``Where's the anchor?'' ``What landmark did
        you read?'' ``What point did you check with?'' ``What would the \textsc{table} show?''
\end{itemize}
\end{skillbox}

\boxguard
\begin{skillbox}[Group Work \& Differentiation]{redbox}
\begin{multicols}{3}
\textbf{Tier R --- Remediate}
\begin{itemize}[leftmargin=1em, itemsep=1pt, topsep=2pt]
  \item $y=(x+1)^2-4$: anchor point, table, plot on a pre-drawn grid.
  \item Same shape, lowest point $(4,0)$ --- write it, then check at $x=6$.
  \item Which of the domain and range did the $-4$ move?
\end{itemize}

\columnbreak
\textbf{Tier A --- Approaching Proficiency}
\begin{itemize}[leftmargin=1em, itemsep=1pt, topsep=2pt]
  \item A parabola with a stretch: vertex and $(4,1)$ give $a=3$.
  \item Verify it at $x=5$ against the graph.
  \item An exponential read from its asymptote: $y=2^x+2$, range $(2,\infty)$.
  \item A gym line from two points: $C(x)=15x+40$, then $C(7)$.
\end{itemize}

\columnbreak
\textbf{Tier E --- Extension}
\begin{itemize}[leftmargin=1em, itemsep=1pt, topsep=2pt]
  \item Find two typos: \texttt{2\^{}X-5} and \texttt{(-X)\^{}2}.
  \item An empty screen: explain it from the anchor, then fix the window from the equation.
  \item Justify: a matching screen proves nothing --- $(x-3)^2$ against $(x-3)^2+0.5$.
\end{itemize}
\end{multicols}
\end{skillbox}

\boxguard[16]
\begin{skillbox}[Individual Work \& Assessment]{redbox}
\textbf{Exit ticket (3 items, no notes):} (1) $y=(x+2)^2-1$ --- the anchor point, the table, and five
points plotted on a pre-drawn grid (\texttt{AF.1f}); (2) an exponential given only by its asymptote
$y=-3$ and the point $(0,-2)$, so the equation must come from the two landmarks, plus the range
(\texttt{AF.1d}); (3) \textbf{the conceptual check} --- a classmate typed \texttt{2\^{}X-3} for ``moved
right $3$,'' which is \emph{the same expression item 2 just produced for a legitimate reason}, and the
student must settle it with one substitution and give the correct keystrokes.

\textbf{Using the results:} item 3 is the one to read, and the pairing with item 2 is deliberate ---
the same keystrokes are right for one job and wrong for another, so no answer can be reached by
pattern-matching. Sort into three piles: \emph{substituted a number}, \emph{restated the rule}, and
\emph{answered from the picture}. The second pile is the biggest and goes to Tier R in 2.4. A student
who writes ``because the $3$ is outside the parentheses'' has repeated 2.2 correctly and answered
nothing today --- the item asks for an input and an output.
\end{skillbox}

\boxguard[18]
\begin{skillbox}[Reinforcement \& Extension]{goldbox}
\textbf{Homework} (10 items): five fluency items --- anchor points, one equation graphed on a
pre-drawn grid, three equations written from verbal descriptions (one per family), one $a$ to solve
for, and a \textbf{keystroke matching table} covering all four of the day's typing cases --- then a
four-part read of \textbf{the yearbook's profit curve}, $P(x)=-2(x-30)^2+1800$, where the staff needs
the profit at \$35 and neither \$35 nor \$1750 sits on a gridline. That is the lesson's argument for
why an equation is worth having.

\textbf{Item 10 is the bridge to Lesson 2.4 and needs no new equation.} It asks the same yearbook
graph for its two zeros, its absolute maximum \emph{with location}, and the two stretches of prices
where profit rises and falls --- all in plain English. Accept plain English; \emph{zero},
\emph{absolute maximum}, and \emph{increasing interval} are 2.4's words to introduce. Insisting on the
vocabulary today buys nothing and costs the bridge.

\textbf{Extension (optional):} $y=-2(x-4)^2+8$ --- all three moves at once, the two inputs where
$y=0$ solved for ($2$ and $6$), and a window that shows them. Then part (c) changes $-2$ to
$-\tfrac12$ and nothing else: the vertex does not move, the parabola flattens, and the two inputs
spread to $0$ and $8$. \textbf{One change to $a$, two visible consequences and one non-consequence}
--- that is the cleanest evidence in the lesson that $a$ controls shape and nothing else.

\textbf{Preview of Lesson 2.4 (\texttt{AFDA.AF.2b--d}):} intercepts, zeros, absolute maximum and
minimum, and the intervals where a graph increases, decreases, or stays constant --- all read from a
graph.
\end{skillbox}

% ── Teacher notes — one per component, in packet order ────────────────────────
% Teacher-only prose lives HERE, never in a _key: a note is the one block in a
% key with no counterpart in the blank, so it makes the key longer than the
% blank. See references/conventions.md ("teachernote").
\boxguard[18]
\begin{teachernote}[Warm-Up]
    \textbf{8 minutes, then score it together.} Answers: (1) $4,1,0,1,4$, and $x=2$ gives $0$;
    (2) $2^5-5=27$ and $2^{5-5}=2^0=1$; (3) (a) parent $x^2$, left $4$; (b) parent $2^x$, down $6$;
    (c) parent $x^2$, reflected over the $x$-axis; (4) domain $(-\infty,\infty)$, range
    $[3,\infty)$; (5) the parent parabola shifted $2$ left, lowest point $(-2,0)$. \textbf{Item 2 is
    the hook, so score it and stop.} Confirm the two numbers, do not explain them, and do not let
    anyone announce which typing is ``right'' --- the hook is worthless if it is resolved here. Item
    1 is not filler: those five outputs are the same arithmetic notes box 1 asks for, done once
    already with no transformation language attached, and a student who is slow here will be slow at
    the plotting. Item 3(c) is the negative-key trap in advance; note who writes $9$ for $-3^2$.
    Item 5 is homework accountability --- if the room is silent, box 1 needs the minutes you
    budgeted for box 3.
\end{teachernote}

\boxguard[18]
\begin{teachernote}[Guided Notes]
    \textbf{Box 1 must include actual plotting; do not narrate it.} The single most common way this
    lesson fails is the teacher filling the image table on the board and then drawing the curve too.
    Fill the table together, then stop talking and let them plot five points --- it takes ninety
    seconds and it is the only place \texttt{AF.1f} is genuinely practised in the notes. The grid is
    pre-drawn with the parent on it, so nobody is sketching from nothing. \textbf{In box 2, the
    exponential row is the hard one and the vocabulary is the reason:} students have been trained for
    a week to find a turning point, and an exponential has none. Point at the dashed line, ask ``where
    was that line before it moved?'' and let them answer $y=0$. The $a$ subsection needs one
    insistence: \emph{write the equation with $a$ in it, then substitute}. Students who substitute
    first produce $a=3$ by coincidence and cannot repeat it. \textbf{Read the warning about lines
    aloud} --- $2f(x)$ and $f(2x)$ both being $y=2x$ is not a technicality, it is why the SOL guidance
    gives lines their own rule, and saying so buys you credibility for the rest of the unit. Box 3 is
    best with a calculator on the projector; if you have one, type the hook's two strings live. If
    time runs short, assign box 3's trap list as reading, but never cut the plotting in box 1.
\end{teachernote}

\boxguard[18]
\begin{teachernote}[Group Activity]
    \textbf{Groups of three, 15 minutes, all groups start at Tier R.} Tier R is one equation graphed
    and one written, and a group that clears it in four minutes should move straight on. Expect Tier R
    item 2 to produce five correct table values and a curve drawn through only three of them --- ask
    for all five plotted before they draw anything. \textbf{Tier A item 1 is where groups reliably go
    wrong}: they use the vertex \emph{as} the second point, get $a=1$, and write the parent's shape
    over a graph that visibly is not. The fix is a question, not a correction: \emph{which two
    different points did you use?} Tier A item 3's range is the parenthesis check ($2$ is never
    reached) --- 2.1's bracket rule paying its fourth dividend. Tier A item 4(c) asks them to justify
    the two-point rule for lines in their own words; a group that gets there has understood something
    most Algebra 2 students have not. \textbf{Tier E item 3 is the share-out that matters.} Read the
    claim aloud --- ``my screen looks exactly like the printed graph, so my equation is right'' --- and
    let the class argue before the group answers. Half of them will defend it, which is exactly why
    the item exists. A group that can say ``half a unit is one pixel'' has the idea; push them to
    finish with the substitution at $x=3$, because naming the flaw is not the same as knowing the
    remedy.
\end{teachernote}

\boxguard[18]
\begin{teachernote}[Exit Ticket]
    \textbf{5 minutes, notes closed, hand it to me at the door.} Answers: (1) anchor $(-2,-1)$, table
    $3,0,-1,0,3$, five points plotted; (2) $y=2^x-3$ from the asymptote at $y=-3$ and the point
    $(0,-2)$, with range $(-3,\infty)$ --- a parenthesis; (3) no: right $3$ requires the image at
    $x=3$ to equal the parent's $2^0=1$, but $2^3-3=5$, so \texttt{2\^{}X-3} is a slide \emph{down}
    $3$ and the correct keystrokes are \texttt{2\^{}(X-3)}. \textbf{Items 2 and 3 are the same
    expression on purpose.} Item 2 produces $2^x-3$ as a correct answer; item 3 shows the identical
    keystrokes failing a different instruction. A student cannot pattern-match their way through
    both, and one who marks item 3 ``correct'' because it matches what they just wrote has told you
    exactly what they are doing. Grade item 3 on whether a \emph{number} appears: ``the $3$ is
    outside the parentheses'' is 2.2's answer, not today's. Do not grade this for points --- use it
    to build the 2.4 groups, where every characteristic gets read off a graph and a student who
    answers from the picture rather than the values will be wrong quietly.
\end{teachernote}

\boxguard[18]
\begin{teachernote}[Homework]
    \textbf{Expect 25--30 minutes.} Item 5 (the keystroke table) is the fastest thing to grade and the
    most diagnostic: the answers run (d), (a), (b), (c), and a student who mismatches (c) and (d) has
    the negative-key confusion that will cost them a whole graph on the test. Item 3(b) is the
    exponential written from an asymptote with no picture at all --- the hardest fluency item here,
    and the one to reteach if it fails. Item 4 is the only place $a$ is solved for without a graph;
    check that the equation is written with $a$ in it before the substitution. \textbf{The yearbook
    items (6--9) are the argument for the whole lesson}, and item 9 is where to spend your grading
    attention: it wants a reason the graph is insufficient (\$35 and \$1750 both fall between
    gridlines), the keystrokes, a window, the \textsc{table} value $1750$, and the range answer.
    Partial credit is fine; ``you can't tell'' with no mention of gridlines is not. \textbf{Item 10 is
    the 2.4 bridge and must be graded on the mathematics, not the vocabulary} --- ``it goes up until
    \$30 then comes down'' is a complete answer today. Insist only on 2.0's rule in 10(b): an extreme
    is a value \emph{and} a location, so \$1800 alone is half an answer. The extension is genuinely
    optional; its part (c) is the best item in the set, because changing $a$ alone moves the two zeros
    and leaves the vertex exactly where it was.
\end{teachernote}
\end{document}
